% ---- score2 standard preamble (from: python3 score2_tools.py preamble) ----
\documentclass[10pt]{extarticle}
\usepackage[a4paper,margin=1.5cm]{geometry}
\usepackage[T1]{fontenc}
\usepackage{amsmath,amssymb}
\usepackage{graphicx}
\usepackage[export]{adjustbox}
\usepackage{multicol}
\usepackage{enumitem}
\usepackage{booktabs,array,tabularx}
\setlength{\columnsep}{0.7cm}
\setlength{\parindent}{0pt}
\setlength{\parskip}{3pt}
\pagestyle{empty}
\makeatletter
% figure registry (figdefs.tex is generated by the extract step)
\newcommand{\qfigdef}[4]{\@namedef{qf@file@#1}{#2}\@namedef{qf@w@#1}{#3}\@namedef{qf@wide@#1}{#4}}
\newcommand{\qf@check}[1]{\@ifundefined{qf@file@#1}{\PackageError{score2}{Unknown figure id #1}{}}{}}
\newcommand{\qf@img}[4][]{\includegraphics[width=#2\dimexpr\@nameuse{qf@w@#4}pt\relax,max width=#3,max height=0.45\textheight,#1]{\@nameuse{qf@file@#4}}}
% Full-width questions: finalize wraps every question that contains a wide figure in \qfullbegin...\qfullend,
% so the WHOLE question (stem, figures, options) is set across the page in reading order, then the two columns
% resume. The question's head - its number, stem and first figure (\qfullsplit marks the end of the head) -
% is kept on one page, so a number is never left at a page foot with its figure overleaf; the rest may break
% across pages normally, so no big gaps are left. Without a \qfullsplit, the whole question is kept together
% if it is under 45% of a page. Never write these yourself.
\newif\ifqf@full
\newif\ifqf@boxed
\newsavebox{\qf@box}
\newcommand{\qf@open}{\global\qf@fulltrue\par\global\qf@boxedtrue
  \setbox\qf@box\vbox\bgroup\hsize\textwidth\linewidth\textwidth}
\newcommand{\qf@emit}[1]{\par\egroup\global\qf@boxedfalse
  \ifdim\dimexpr\ht\qf@box+\dp\qf@box\relax<#1\textheight\noindent\usebox{\qf@box}\par\else\unvbox\qf@box\fi}
\newcommand{\qfullsplit}{\ifqf@boxed\qf@emit{0.85}\fi}
\newcommand{\qf@close}{\ifqf@boxed\qf@emit{0.45}\fi\global\qf@fullfalse}
\newcommand{\qfullbegin}{\end{multicols}\qf@open}
\newcommand{\qfullend}{\qf@close\begin{multicols}{2}}
\newcommand{\qfullstart}{\qf@open}
\newcommand{\qfulldone}{\qf@close}
% \qfig[scale]{id}: block figure, centred, directly where it is written (size set by the tool).
% Fallback if a wide figure is not inside a full-width question: break out of the columns in place - no floats.
\newcommand{\qfig}[2][1]{\qf@check{#2}%
  \ifqf@full\qf@blk{#1}{#2}\else
  \ifnum\@nameuse{qf@wide@#2}=1 \expandafter\@firstoftwo\else\expandafter\@secondoftwo\fi
  {\end{multicols}\begin{center}\qf@img{#1}{\textwidth}{#2}\end{center}\begin{multicols}{2}}%
  {\qf@blk{#1}{#2}}\fi}
\newcommand{\qf@blk}[2]{\par\smallskip{\centering\qf@img{#1}{\linewidth}{#2}\par}\smallskip}
% \qfigi[scale]{id}: inline figure (inside option lists / tables / minipages); never breaks columns.
\newcommand{\qfigi}[2][1]{\qf@check{#2}\qf@img[valign=c]{#1}{\linewidth}{#2}}
\makeatother
% structure macros - lint reads these to check completeness and figure placement
\newcommand{\papersec}[1]{\par\bigskip{\centering\Large\bfseries PAPER-#1\par}\medskip}
\newcommand{\partsec}[2]{\par\medskip\noindent{\large\bfseries PART-#1: #2}\par\smallskip}
\newcommand{\sect}[1]{\par\smallskip\noindent\textbf{#1}\par}
\newcommand{\qn}[1]{\par\smallskip\noindent\textbf{#1.}~}
\newcommand{\paraq}[2]{\par\smallskip\noindent\textbf{Paragraph for Questions #1 and #2}\par}
\newenvironment{qoptions}{\begin{itemize}[leftmargin=1.8em,itemsep=1pt,topsep=2pt,label={}]}{\end{itemize}}
\newcommand{\opt}[2]{\item[(#1)] #2}
% figure registry (generated by score2_tools.py finalize)
\qfigdef{p1_phy_s1_q01_1}{Score-I - Test 12 (09 Mar 2022)-img/p1_phy_s1_q01_1.png}{129.4}{0}
\qfigdef{p1_phy_s1_q02_1}{Score-I - Test 12 (09 Mar 2022)-img/p1_phy_s1_q02_1.png}{103.2}{0}
\qfigdef{p1_phy_s1_q03_1}{Score-I - Test 12 (09 Mar 2022)-img/p1_phy_s1_q03_1.png}{92.0}{0}
\qfigdef{p1_phy_s1_q04_1}{Score-I - Test 12 (09 Mar 2022)-img/p1_phy_s1_q04_1.png}{135.6}{0}
\qfigdef{p1_phy_s1_q05_1}{Score-I - Test 12 (09 Mar 2022)-img/p1_phy_s1_q05_1.png}{126.6}{0}
\qfigdef{p1_phy_s1_q06_1}{Score-I - Test 12 (09 Mar 2022)-img/p1_phy_s1_q06_1.png}{98.0}{0}
\qfigdef{p1_phy_s1_q07_1}{Score-I - Test 12 (09 Mar 2022)-img/p1_phy_s1_q07_1.png}{137.3}{0}
\qfigdef{p1_phy_s1_q08_1}{Score-I - Test 12 (09 Mar 2022)-img/p1_phy_s1_q08_1.png}{131.7}{0}
\qfigdef{p1_phy_s1_q10_1}{Score-I - Test 12 (09 Mar 2022)-img/p1_phy_s1_q10_1.png}{97.5}{0}
\qfigdef{p1_phy_s2_q01_1}{Score-I - Test 12 (09 Mar 2022)-img/p1_phy_s2_q01_1.png}{128.9}{0}
\qfigdef{p1_phy_s2_q03_1}{Score-I - Test 12 (09 Mar 2022)-img/p1_phy_s2_q03_1.png}{187.3}{0}
\qfigdef{p1_phy_s3_q07_1}{Score-I - Test 12 (09 Mar 2022)-img/p1_phy_s3_q07_1.png}{92.3}{0}
\qfigdef{p1_phy_s3_q08_1}{Score-I - Test 12 (09 Mar 2022)-img/p1_phy_s3_q08_1.png}{130.8}{0}
\qfigdef{p1_phy_s3_q09_1}{Score-I - Test 12 (09 Mar 2022)-img/p1_phy_s3_q09_1.png}{129.3}{0}
\qfigdef{p1_phy_s3_q10_1}{Score-I - Test 12 (09 Mar 2022)-img/p1_phy_s3_q10_1.png}{85.5}{0}
\qfigdef{p1_chem_s1_q01_1}{Score-I - Test 12 (09 Mar 2022)-img/p1_chem_s1_q01_1.png}{127.6}{0}
\qfigdef{p1_chem_s1_q01_A}{Score-I - Test 12 (09 Mar 2022)-img/p1_chem_s1_q01_A.png}{55.4}{0}
\qfigdef{p1_chem_s1_q01_C}{Score-I - Test 12 (09 Mar 2022)-img/p1_chem_s1_q01_C.png}{58.2}{0}
\qfigdef{p1_chem_s1_q01_D}{Score-I - Test 12 (09 Mar 2022)-img/p1_chem_s1_q01_D.png}{51.2}{0}
\qfigdef{p1_chem_s1_q02_A}{Score-I - Test 12 (09 Mar 2022)-img/p1_chem_s1_q02_A.png}{158.4}{0}
\qfigdef{p1_chem_s1_q02_B}{Score-I - Test 12 (09 Mar 2022)-img/p1_chem_s1_q02_B.png}{134.9}{0}
\qfigdef{p1_chem_s1_q02_C}{Score-I - Test 12 (09 Mar 2022)-img/p1_chem_s1_q02_C.png}{136.4}{0}
\qfigdef{p1_chem_s1_q02_4}{Score-I - Test 12 (09 Mar 2022)-img/p1_chem_s1_q02_4.png}{55.4}{0}
\qfigdef{p1_chem_s1_q03_B}{Score-I - Test 12 (09 Mar 2022)-img/p1_chem_s1_q03_B.png}{60.1}{0}
\qfigdef{p1_chem_s1_q03_C}{Score-I - Test 12 (09 Mar 2022)-img/p1_chem_s1_q03_C.png}{67.7}{0}
\qfigdef{p1_chem_s1_q03_D}{Score-I - Test 12 (09 Mar 2022)-img/p1_chem_s1_q03_D.png}{68.1}{0}
\qfigdef{p1_chem_s1_q04_1}{Score-I - Test 12 (09 Mar 2022)-img/p1_chem_s1_q04_1.png}{150.6}{0}
\qfigdef{p1_chem_s2_q01_1}{Score-I - Test 12 (09 Mar 2022)-img/p1_chem_s2_q01_1.png}{311.5}{1}
\qfigdef{p1_chem_s2_q02_1}{Score-I - Test 12 (09 Mar 2022)-img/p1_chem_s2_q02_1.png}{71.1}{0}
\qfigdef{p1_chem_s3_q06_1}{Score-I - Test 12 (09 Mar 2022)-img/p1_chem_s3_q06_1.png}{164.8}{0}
\qfigdef{p1_chem_s3_q07_1}{Score-I - Test 12 (09 Mar 2022)-img/p1_chem_s3_q07_1.png}{220.1}{0}
\qfigdef{p2_phy_s1_para3_4_1}{Score-I - Test 12 (09 Mar 2022)-img/p2_phy_s1_para3_4_1.png}{66.1}{0}
\qfigdef{p2_phy_s1_para3_4_2}{Score-I - Test 12 (09 Mar 2022)-img/p2_phy_s1_para3_4_2.png}{77.0}{0}
\qfigdef{p2_phy_s1_para3_4_3}{Score-I - Test 12 (09 Mar 2022)-img/p2_phy_s1_para3_4_3.png}{67.5}{0}
\qfigdef{p2_phy_s1_para3_4_4}{Score-I - Test 12 (09 Mar 2022)-img/p2_phy_s1_para3_4_4.png}{52.9}{0}
\qfigdef{p2_phy_s1_para5_6_1}{Score-I - Test 12 (09 Mar 2022)-img/p2_phy_s1_para5_6_1.png}{112.5}{0}
\qfigdef{p2_phy_s1_para7_8_1}{Score-I - Test 12 (09 Mar 2022)-img/p2_phy_s1_para7_8_1.png}{61.7}{0}
\qfigdef{p2_phy_s1_para7_8_5}{Score-I - Test 12 (09 Mar 2022)-img/p2_phy_s1_para7_8_5.png}{70.8}{0}
\qfigdef{p2_phy_s1_para7_8_9}{Score-I - Test 12 (09 Mar 2022)-img/p2_phy_s1_para7_8_9.png}{56.8}{0}
\qfigdef{p2_phy_s1_para7_8_2}{Score-I - Test 12 (09 Mar 2022)-img/p2_phy_s1_para7_8_2.png}{63.9}{0}
\qfigdef{p2_phy_s1_para7_8_6}{Score-I - Test 12 (09 Mar 2022)-img/p2_phy_s1_para7_8_6.png}{81.7}{0}
\qfigdef{p2_phy_s1_para7_8_10}{Score-I - Test 12 (09 Mar 2022)-img/p2_phy_s1_para7_8_10.png}{97.3}{0}
\qfigdef{p2_phy_s1_para7_8_3}{Score-I - Test 12 (09 Mar 2022)-img/p2_phy_s1_para7_8_3.png}{51.9}{0}
\qfigdef{p2_phy_s1_para7_8_7}{Score-I - Test 12 (09 Mar 2022)-img/p2_phy_s1_para7_8_7.png}{89.3}{0}
\qfigdef{p2_phy_s1_para7_8_11}{Score-I - Test 12 (09 Mar 2022)-img/p2_phy_s1_para7_8_11.png}{93.0}{0}
\qfigdef{p2_phy_s1_para7_8_4}{Score-I - Test 12 (09 Mar 2022)-img/p2_phy_s1_para7_8_4.png}{57.0}{0}
\qfigdef{p2_phy_s1_para7_8_8}{Score-I - Test 12 (09 Mar 2022)-img/p2_phy_s1_para7_8_8.png}{70.5}{0}
\qfigdef{p2_phy_s1_para7_8_12}{Score-I - Test 12 (09 Mar 2022)-img/p2_phy_s1_para7_8_12.png}{93.5}{0}
\qfigdef{p2_phy_s2_q01_1}{Score-I - Test 12 (09 Mar 2022)-img/p2_phy_s2_q01_1.png}{158.2}{0}
\qfigdef{p2_phy_s2_q03_1}{Score-I - Test 12 (09 Mar 2022)-img/p2_phy_s2_q03_1.png}{150.4}{0}
\qfigdef{p2_phy_s2_q05_1}{Score-I - Test 12 (09 Mar 2022)-img/p2_phy_s2_q05_1.png}{107.2}{0}
\qfigdef{p2_phy_s2_q06_1}{Score-I - Test 12 (09 Mar 2022)-img/p2_phy_s2_q06_1.png}{108.9}{0}
\qfigdef{p2_phy_s3_q07_1}{Score-I - Test 12 (09 Mar 2022)-img/p2_phy_s3_q07_1.png}{6.3}{0}
\qfigdef{p2_phy_s3_q07_2}{Score-I - Test 12 (09 Mar 2022)-img/p2_phy_s3_q07_2.png}{20.8}{0}
\qfigdef{p2_phy_s3_q07_3}{Score-I - Test 12 (09 Mar 2022)-img/p2_phy_s3_q07_3.png}{21.5}{0}
\qfigdef{p2_phy_s3_q08_1}{Score-I - Test 12 (09 Mar 2022)-img/p2_phy_s3_q08_1.png}{153.9}{0}
\qfigdef{p2_phy_s3_q09_1}{Score-I - Test 12 (09 Mar 2022)-img/p2_phy_s3_q09_1.png}{89.4}{0}
\qfigdef{p2_phy_s3_q09_2}{Score-I - Test 12 (09 Mar 2022)-img/p2_phy_s3_q09_2.png}{115.5}{0}
\qfigdef{p2_phy_s3_q09_3}{Score-I - Test 12 (09 Mar 2022)-img/p2_phy_s3_q09_3.png}{95.0}{0}
\qfigdef{p2_phy_s3_q10_1}{Score-I - Test 12 (09 Mar 2022)-img/p2_phy_s3_q10_1.png}{63.4}{0}
\qfigdef{p2_phy_s3_q10_2}{Score-I - Test 12 (09 Mar 2022)-img/p2_phy_s3_q10_2.png}{67.1}{0}
\qfigdef{p2_phy_s3_q10_3}{Score-I - Test 12 (09 Mar 2022)-img/p2_phy_s3_q10_3.png}{67.1}{0}
\qfigdef{p2_chem_s1_para1_2_1}{Score-I - Test 12 (09 Mar 2022)-img/p2_chem_s1_para1_2_1.png}{65.6}{0}
\qfigdef{p2_chem_s1_para1_2_2}{Score-I - Test 12 (09 Mar 2022)-img/p2_chem_s1_para1_2_2.png}{107.0}{0}
\qfigdef{p2_chem_s1_para1_2_3}{Score-I - Test 12 (09 Mar 2022)-img/p2_chem_s1_para1_2_3.png}{107.0}{0}
\qfigdef{p2_chem_s1_para1_2_4}{Score-I - Test 12 (09 Mar 2022)-img/p2_chem_s1_para1_2_4.png}{57.3}{0}
% ---- body skeleton ----
% \begin{document}
% \begin{center}{\Large\bfseries JEE (Main+Advanced) Score-II --- Test NN (DD Mon YYYY)}\end{center}
% \begin{multicols}{2}
%   \papersec{1}   (only for two-paper tests)
%   \partsec{1}{PHYSICS} \sect{Section-I (i): Single Correct} \qn{1} ... \qfig{p1_phy_q01_1} ...
% \end{multicols}
% \end{document}

\begin{document}
\begin{center}{\Large\bfseries JEE (Main+Advanced) Score-I --- Test 12 (09 Mar 2022)}\end{center}
\raggedcolumns
\setlength{\tabcolsep}{3pt}
\begin{multicols}{2}
\papersec{1}
\partsec{1}{PHYSICS}
\sect{SECTION-I}
\qn{1} The circuit diagram shows a source of emf 10 V having an internal resistance of 2.5\,$\Omega$ connected to an ideal ammeter. The jockey can be slided to different points on the Rheostat whose total resistance is 10\,$\Omega$. What can be the reading of the ammeter?
\qfig{p1_phy_s1_q01_1}
\begin{qoptions}
\opt{A}{1 A}
\opt{B}{2 A}
\opt{C}{3 A}
\opt{D}{4 A}
\end{qoptions}

\qn{2} Consider the LC circuit below. At the time shown the current has its maximum value. At this time which of the following is/are incorrect?
\qfig{p1_phy_s1_q02_1}
\begin{qoptions}
\opt{A}{The charge on the capacitor has its maximum value}
\opt{B}{The magnetic field inside the inductor is zero}
\opt{C}{The electric field inside the capacitor has its maximum value}
\opt{D}{The charge on the capacitor is zero}
\end{qoptions}

\qn{3} A particle of mass $m$ and charge $-q$ has been projected from ground as shown in the figure below. Mark out the correct statement(s):
\qfig{p1_phy_s1_q03_1}
\begin{qoptions}
\opt{A}{The path of motion of the particle is parabolic}
\opt{B}{The path of motion of the particle is a straight line}
\opt{C}{Time of flight of particle is $\dfrac{2u\sin\theta}{g}$}
\opt{D}{Range of motion of the particle will be less than $\dfrac{u^2\sin2\theta}{g}$}
\end{qoptions}

\qn{4} Consider a situation in which there is a Newtonian fluid between two long coaxial uniformly moving cylinders of radius $R$ and $2R$.
\qfig{p1_phy_s1_q04_1}
\begin{qoptions}
\opt{A}{Force per unit length required to keep the inner cylinder moving is $4\pi\eta v_0$}
\opt{B}{Force per unit length required to keep the inner cylinder moving is $\dfrac{6\pi\eta v_0}{\ell n2}$}
\opt{C}{Distance of the point from the axis at which speed of fluid is zero is $2^{1/3}R$}
\opt{D}{Distance of the point from the axis at which speed of fluid is zero is $\sqrt2\,R$}
\end{qoptions}

\qn{5} In the figure shown the radius of curvature of the left \& right surface of the concave lens are 10 cm \& 15 cm respectively. The radius of curvature of the mirror is 15 cm:
\qfig{p1_phy_s1_q05_1}
\begin{qoptions}
\opt{A}{modulus of equivalent focal length of the combination is 18 cm}
\opt{B}{modulus of equivalent focal length of the combination is 36 cm}
\opt{C}{the system behaves like a concave mirror}
\opt{D}{the system behaves like a convex mirror}
\end{qoptions}

\qn{6} A spool of mass $m$ and moment of inertia $I$ about its axis rests on a smooth inclined plane, inclined at an angle $\theta$. The end of the thread is tied as shown. The outer radius of the spool is $R$ \& inner radius of the spool is $r$:
\qfig{p1_phy_s1_q06_1}
\begin{qoptions}
\opt{A}{Acceleration of centre of mass of spool $a_{cm}=\dfrac{mg\sin\theta}{m+\frac{I}{r^2}}$}
\opt{B}{Tension in thread is $T=\dfrac{mgI\sin\theta}{I+mr^2}$}
\opt{C}{Angular acceleration of spool is $\alpha=\dfrac{a_{cm}}{R}$}
\opt{D}{Acceleration of centre of mass is zero.}
\end{qoptions}

\qn{7} A particle executes SHM about a point other than $x=0$ as shown in the graph. CHOOSE the correct option(s):
\qfig{p1_phy_s1_q07_1}
\begin{qoptions}
\opt{A}{Amplitude is equal to 4 m}
\opt{B}{Angular frequency $=\dfrac\pi3$}
\opt{C}{Equilibrium position is at $x=2$ m}
\opt{D}{Angular frequency $=\dfrac{2\pi}3$}
\end{qoptions}

\qn{8} $M_1$ and $M_2$ are two plane mirror \& $S$ is monochromatic source. $AB$ is a `stop' which stops direct light to reach $CF$ and allows reflected light from $M_1$ and $M_2$ to reach $CF$. $P$ is a point on $CF$ such that $SP$ is parallel to $M_1$ \& $M_2$ and perpendicular to $CF$:
\qfig{p1_phy_s1_q08_1}
\begin{qoptions}
\opt{A}{Circular fringes will be formed on $CF$}
\opt{B}{$P$ will be central maxima}
\opt{C}{$P$ will be a point of maxima if $\lambda=\dfrac{6d^2}{D}$}
\opt{D}{$P$ will be a point of maxima if $\lambda=\dfrac{d^2}{D}$}
\end{qoptions}

\qn{9} Two vehicles, each moving with speed $u$ on the same horizontal straight road, are approaching each other. Wind blows along the road with velocity $w$. One of these vehicles blows a whistle of frequency $f_1$. An observer in the other vehicle hears the frequency of the whistle to be $f_2$. The speed of sound in still air is $V$. The correct statement(s) is (are):
\begin{qoptions}
\opt{A}{If the wind blows from the observer to the source, $f_2>f_1$}
\opt{B}{If the wind blows from the source to the observer, $f_2>f_1$}
\opt{C}{If the wind blows from observer to the source, $f_2<f_1$}
\opt{D}{If the wind blows from the source to the observer, $f_2<f_1$}
\end{qoptions}

\qn{10} Two moles of helium gas are taken over the cycle $ABCDA$, as shown in the $P$--$T$ diagram. (Use: $\ln2=0.7$)
\qfig{p1_phy_s1_q10_1}
\begin{qoptions}
\opt{A}{The work done by the gas from $A$ to $B$ is $600R$}
\opt{B}{The work done by the gas from $B$ to $C$ is $840R$}
\opt{C}{The net work done by the gas in the cycle $ABCDA$ is $276R$}
\opt{D}{The net work done on the gas in the cycle $ABCDA$ is $1076R$}
\end{qoptions}

\sect{SECTION-II(i)}
\qn{1} An infinitely long straight wire carries a current $i_1=1$ amp. A rectangular current loop placed in the same plane of shown dimensions carries a current $i_2=1$ amp anticlockwise. Magnitude of net force on the loop is found to be $\beta\dfrac{\mu_0}{4\pi}$. Find value of $\beta$.
\qfig{p1_phy_s2_q01_1}

\qn{2} A single electron orbits a stationary nucleus of charge $+Ze$, where $Z$ is a constant and $e$ is the magnitude of the electronic charge. It requires 47.2 eV to excite the electron from the second Bohr orbit to third Bohr orbit. Find the value of $Z$.

\qn{3} The positions of ${}^2_1$D, ${}^4_2$He and ${}^7_3$Li are shown on the binding energy per nucleon curve as shown in figure. Find the energy released in the fusion reaction, ${}^2_1\text{D}+{}^7_3\text{Li}\to2\,{}^4_2\text{He}+{}^1_0\text{n}$ in multiple of 5.5 MeV.
\qfig{p1_phy_s2_q03_1}

\qn{4} A string fixed at both the ends of length 2 m, vibrating in its 7th overtone. Equation of the standing wave is given by $y=A\sin kx\cos(\omega t+\pi/3)$. All the symbols have their usual meaning. Mass per unit length of the string is 0.5 gm/cm. Given that $A=1$ cm and $\omega=100\pi$ rad/sec. If after time $t_0$ starting from $t=0$ sec, energy of vibration is completely potential for the second time. Find $120\,t_0$. (Use $\pi^2=10$)

\qn{5} A particle of mass 2 kg moves on a smooth horizontal plane under the action of a single force $\vec F=(3\hat i+4\hat j)$ N. Under this force it is displaced from $(0,0)$ to $(1\text{ m},1\text{ m})$. The initial speed of the particle is $\sqrt3$ m/s and final speed is $\sqrt{5k}$ m/s. Find the value of $k$.

\sect{SECTION-II(ii)}
\qn{6} The height at which the acceleration due to gravity decreases by 36\% of its value on the surface of the earth is $R/x$, then the value of $x$ is: (Assume radius of the earth is $R$)

\qn{7} Find the charge in $\mu$C on $10\,\mu$F capacitor in steady state. If the value of charge is $x\,\mu$C, then find the value of $x/100$.
\qfig{p1_phy_s3_q07_1}

\qn{8} The velocities of two steel balls before impact are shown. After head on elastic collision find the velocity (in m/s) of ball $B$.
\qfig{p1_phy_s3_q08_1}

\qn{9} In the circuit indicated, find the current through both the batteries $i_1$ \& $i_2$. If $|i_1|+|i_2|=x$ then find the value of $\dfrac{5x}{6}$.
\qfig{p1_phy_s3_q09_1}

\qn{10} In an L-C circuit shown in the figure, $C=1$ F, $L=4$ H. At time $t=0$, charge in the capacitor is 4 C and it is decreasing at a rate of $\sqrt5$ C/s. The maximum charge in the capacitor is $n$ C, then value of $n$ will be:
\qfig{p1_phy_s3_q10_1}

\partsec{2}{CHEMISTRY}
\sect{SECTION-I}
\qn{1} \qfig{p1_chem_s1_q01_1}
\begin{center}\small
\begin{tabular}{@{}l@{\ }c@{\qquad\qquad}l@{\ }c@{}}
(A) & \qfigi{p1_chem_s1_q01_A}\rule{0pt}{40pt} & (B) & CH$_3$--CHO\\[8pt]
(C) & \qfigi{p1_chem_s1_q01_C}\rule{0pt}{26pt} & (D) & \qfigi{p1_chem_s1_q01_D}
\end{tabular}
\end{center}

\qn{2} Which reactions are correctly matched with their product:
\begin{qoptions}
\opt{A}{\qfigi{p1_chem_s1_q02_A}}
\opt{B}{\qfigi{p1_chem_s1_q02_B}}
\opt{C}{\qfigi{p1_chem_s1_q02_C}}
\opt{D}{(CH$_3$COO)$_2$Ca \qfigi{p1_chem_s1_q02_4}}
\end{qoptions}

\qn{3} Vicinal-dichloride is obtained as a product in reaction -
\begin{center}\small
\begin{tabular}{@{}l@{\ }l@{\qquad}l@{\ }c@{}}
(A) & CH$_3$--CH=O $\xrightarrow{\text{PCl}_5}$ & (B) & \qfigi{p1_chem_s1_q03_B}\rule{0pt}{26pt}\\[8pt]
(C) & \qfigi{p1_chem_s1_q03_C}\rule{0pt}{32pt} & (D) & \qfigi{p1_chem_s1_q03_D}\rule{0pt}{40pt}
\end{tabular}
\end{center}

\qn{4} \qfig{p1_chem_s1_q04_1}
Correct option is/are:
\begin{qoptions}
\opt{A}{H$_2$O $>$ P $>$ R \ (Acidic strength)}
\opt{B}{Q $>$ P \ (Rate of dehydration)}
\opt{C}{All P, Q \& R gives same dehydration major product with conc. H$_2$SO$_4$}
\opt{D}{All P, Q \& R react with Lucas reagent with same rate}
\end{qoptions}

\qn{5} Select the correct statement(s) regarding $3p_y$ orbital
\begin{qoptions}
\opt{A}{Total no. of nodes are 2}
\opt{B}{Number of maxima in the curve $4\pi r^2R^2$ vs $r$ are two}
\opt{C}{Quantum no. $n$, $\ell$ and $m$ for orbital may be 3, 1, $-1$ respectively}
\opt{D}{The magnetic quantum number may have a positive value}
\end{qoptions}

\qn{6} Which is/are correct statement(s) about the solubility of AgCl(s).\\ {Given: $K_{sp(\text{AgCl})}=10^{-10}$; $K_f[\text{Ag(NH}_3)_2^+]=10^8$.}
\begin{qoptions}
\opt{A}{Solubility of AgCl in pure water is $10^{-5}$ gm/litre}
\opt{B}{Solubility of AgCl in 2M KBr is $10^{-5}$ gm/litre}
\opt{C}{Solubility of AgCl in 2M AgNO$_3$ is $5\times10^{-11}$ M}
\opt{D}{Solubility of AgCl in 2M NH$_3$ is 0.166 M}
\end{qoptions}

\qn{7} Select the correct statement(s)
\begin{qoptions}
\opt{A}{bombarding ${}_7$N$^{14}$ nuclei with $\alpha$-particle, the nuclei of the product formed after release of proton would be ${}_8$O$^{17}$}
\opt{B}{Decay constant does not depend upon temperature}
\opt{C}{Nuclide and it's decay product after $\alpha$-emission are called isodiaphers}
\opt{D}{Half-life of radium is 1580 years. It's average life will be 1097.22 years.}
\end{qoptions}

\qn{8} In electrofining of metals anode and cathode are taken as thick slab of impure metal and strip of pure metal respectively while the electrolyte is aqueous solution of a complex metal salt. This method cannot be applied for the refining of :
\begin{qoptions}
\opt{A}{Copper}
\opt{B}{Sodium}
\opt{C}{Aluminium}
\opt{D}{Silver}
\end{qoptions}

\qn{9} Which of the following sulphide(s) does/do not liberate H$_2$S on warning with dil. HCl.
\begin{qoptions}
\opt{A}{HgS}
\opt{B}{ZnS}
\opt{C}{FeS}
\opt{D}{CuS}
\end{qoptions}

\qn{10} Consider the following transformation:\\
CuSO$_{4(\text{aq.})}$ + KI$_{(\text{excess})}$ $\to$ Product\\
Select the INCORRECT statement.
\begin{qoptions}
\opt{A}{Product contains [Cu(H$_2$O)$_4$]$^{2+}$ ions.}
\opt{B}{Presence of brown colour in product is due to I$_3^-$ ion.}
\opt{C}{Oxidation state of sulphur in reactant and product is different}
\opt{D}{White ppt. of CuI$_2$ is observed in product}
\end{qoptions}

\qfullbegin
\sect{SECTION-II(i)}
\qn{1} How many compound(s) gives Iodoform test with NaOH + I$_2$:
\qfig{p1_chem_s2_q01_1}
\qfullend

\qn{2} Initially one mole of ideal gas $\left(C_v=\dfrac52R\right)$ at 1.0 atm and 300 K is put through the following cycle:\\
Step-I: Heating to twice its initial pressure at constant volume. ($R=0.082\left(\dfrac{\text{litre-atm}}{\text{K-mole}}\right)$, $\sqrt2=1.4$)\\
Step-II: Adiabatic expansion to its initial temperature.\\
Step-III: Isothermal compression back to 1.00 atm.\\
What is the volume at state $X$ (in litre)?
\qfig{p1_chem_s2_q02_1}

\qn{3} When 1.0 g of oxalic acid (H$_2$C$_2$O$_4$) is burned in a bomb calorimeter whose heat capacity 8.75 kJ/K, the temperature increases by 0.312 K. The enthalpy of combustion of oxalic acid 27$^\circ$C is (kJ/mole) ($R=8.3$ J/mole-K)

\qn{4} RH$_2$ (ion exchange resin) can replace Ca$^{2+}$ ions in hard water as RH$_2$ + Ca$^{2+}\to$ RCa + 2H$^+$. If 8 Litre hard water after passing through RH$_2$ has pH = 5, then hardness in parts per million of Ca$^{2+}$ is.

\qn{5} N$_2$(g) reacts with H$_2$(g) in either of the following ways depending upon \& supply of H$_2$(g).\\
N$_2$(g) + H$_2$(g) $\to$ N$_2$H$_2$(g)\\
N$_2$(g) + 2H$_2$(g) $\to$ N$_2$H$_4$(g)\\
If 3 litre N$_{2(g)}$ and 5 litre H$_{2(g)}$ are taken initially (at same temperature and pressure). Calculate the contraction in volume after the reaction (in L)

\sect{SECTION-II(ii)}
\qn{6} Consider following sequence of reactions
\qfig{p1_chem_s3_q06_1}
Total sum of degree of unsaturations of product D \& E will be

\qn{7} How many of the following compounds will show higher rate of electrophilic substitution reaction than benzene -
\qfig{p1_chem_s3_q07_1}

\qn{8} Find the total number of compounds which on heating undergo redox reactions.\\
PbCl$_4$, Mg(NO$_3$)$_2$, HgC$_2$O$_4$, Ag$_2$CO$_3$, Pb(CN)$_4$, Al(OH)$_3$, Cu(CN)$_2$

\qn{9} Find total number of complex in which metal-carbon linkage(s) is/are present.\\\relax
[Fe(EDTA)]$^-$, K[PtCl$_3$($\pi$-C$_2$H$_4$)], [Cr(acac)$_3$], [Fe(CO)$_5$], [Co(C$_2$H$_5$O)$_6$]$^{3-}$, [Ni(dmg)$_2$]

\qn{10} Consider the following interaction\\
NaF + XeF$_4\to$ P (ionic compound)\\
Find the value of $\left[\dfrac{X-Y}{Z}\right]$\\
where\\
$X$ = Value of adjacent bond angle in XeF$_4$\\
$Y$ = Value of adjacent bond angle in anionic part of compound `P'\\
$Z$ = Maximum number of atoms that can lie in a plane in anionic part of compound `P'

\partsec{3}{MATHEMATICS}
\sect{SECTION-I}
\qn{1} If $\alpha,\beta$ be the roots of the equation $\mu^2-2\mu+2=0$ and if $\cot\theta=x+1$, then $\dfrac{(x+\alpha)^n-(x+\beta)^n}{\alpha-\beta}$ is equal to, {(where $\theta\ne k_1\pi$, $k_1\in I$)}
\begin{qoptions}
\opt{A}{$\dfrac{\sin n\theta}{\sin^n\theta}$}
\opt{B}{$\dfrac{\cos n\theta}{\cos^n\theta}$}
\opt{C}{$\dfrac{\sin n\theta}{\cos^n\theta}$}
\opt{D}{$\dfrac{\operatorname{cosec}^n\theta}{\operatorname{cosec}n\theta}$, where $n\theta\ne m\pi$, $m\in I$}
\end{qoptions}

\qn{2} If equation of tangent at $P$, $Q$ and vertex $A$ of a parabola are $3x+4y-7=0$, $2x+3y-10=0$ and $x-y=0$ respectively, then
\begin{qoptions}
\opt{A}{focus is $(4,5)$}
\opt{B}{length of latus rectum is $2\sqrt2$}
\opt{C}{axis is $x+y-9=0$}
\opt{D}{vertex is $\left(\dfrac92,\dfrac92\right)$}
\end{qoptions}

\qn{3} If foci of $\dfrac{x^2}{a^2}-\dfrac{y^2}{b^2}=1$ concide with the focii of $\dfrac{x^2}{25}+\dfrac{y^2}9=1$ and eccentricity of the hyperbola is 2, then
\begin{qoptions}
\opt{A}{$a^2+b^2=16$}
\opt{B}{there is no director circle to the hyperbola}
\opt{C}{centre of the director circle of hyperbola is $(0,0)$}
\opt{D}{length of latus ractum of the hyperbola $=12$}
\end{qoptions}

\qn{4} If $I_n=\displaystyle\int_0^1\dfrac{dx}{(1+x^2)^n}$, $n\in N$, then which of the following statements hold good?
\begin{qoptions}
\opt{A}{$2nI_{n+1}=2^{-n}+(2n-1)I_n$}
\opt{B}{$I_2=\dfrac\pi8+\dfrac14$}
\opt{C}{$I_2=\dfrac\pi8-\dfrac14$}
\opt{D}{$I_3=\dfrac\pi{16}-\dfrac5{48}$}
\end{qoptions}

\qn{5} If $F(x)=\dfrac1{x^2}\displaystyle\int_4^x\left(4t^2-2F'(t)\right)dt$, then $F'(4)$ equals --
\begin{qoptions}
\opt{A}{$\dfrac{32}9$}
\opt{B}{$\dfrac{64}9$}
\opt{C}{$\dfrac{F(8)}{28}$}
\opt{D}{$\dfrac{11F(8)}{28}$}
\end{qoptions}

\qn{6} The solution of $\left(\dfrac{dy}{dx}\right)(x^2y^3+xy)=1$ is, (where $C$ is arbitrary constant)
\begin{qoptions}
\opt{A}{$1/x=2-y^2+Ce^{-y^2/2}$}
\opt{B}{$2/x=1-y^2+Ce^{-y/2}$}
\opt{C}{$2/x=1+y^2+Ce^{-y/2}$}
\opt{D}{$\dfrac{1-2x}{x}=-y^2+Ce^{-y^2/2}$}
\end{qoptions}

\qn{7} Let $f(x)=(x-1)^4(x-2)^n$, $n\in N$. $f(x)$ has
\begin{qoptions}
\opt{A}{Local minimum at $x=2$ if $n$ is even}
\opt{B}{Local minimum at $x=1$ if $n$ is odd}
\opt{C}{Local maximum at $x=1$ if $n$ is odd}
\opt{D}{Local minimum at $x=1$ if $n$ is even}
\end{qoptions}

\qn{8} Let $A=\begin{bmatrix}1&2&2\\2&1&2\\2&2&1\end{bmatrix}$, then
\begin{qoptions}
\opt{A}{$A^2-4A-5I_3=0$}
\opt{B}{$A^{-1}=\dfrac15(A-4I_3)$}
\opt{C}{$A^3$ is not invertible}
\opt{D}{$A^2$ is invertible}
\end{qoptions}

\qn{9} System of equation\\ $x+3y+2z=6$\\ $x+\lambda y+2z=7$\\ $x+3y+2z=\mu$ has
\begin{qoptions}
\opt{A}{unique solution if $\lambda=2$, $\mu\ne6$}
\opt{B}{infinitely many solution if $\lambda=4$, $\mu=6$}
\opt{C}{no solution if $\lambda=5$, $\mu=7$}
\opt{D}{no solution if $\lambda=3$, $\mu=5$}
\end{qoptions}

\qn{10} The plane $\ell x+my=0$ is rotated about its line of intersection with the plane $z=0$, through an angle $\alpha$, then equation of plane in its new position may be
\begin{qoptions}
\opt{A}{$\ell x+my+z\sqrt{\ell^2+m^2}\tan\alpha=0$}
\opt{B}{$\ell x+my-z\sqrt{\ell^2+m^2}\tan\alpha=0$}
\opt{C}{data is not sufficient}
\opt{D}{None of these}
\end{qoptions}

\sect{SECTION-II(i)}
\qn{1} The values of `a' for which the function $f(x)=\sin x-a\sin2x-\dfrac13\sin3x+2ax$ increases throughout the number line is $[\lambda,\infty)$, then find $\dfrac{100\lambda}3$

\qn{2} If $f(x)=\begin{vmatrix}\cos(x+\alpha)&\cos(x+\beta)&\cos(x+\gamma)\\\sin(x+\alpha)&\sin(x+\beta)&\sin(x+\gamma)\\\sin(\beta-\gamma)&\sin(\gamma-\alpha)&\sin(\alpha-\beta)\end{vmatrix}$ and $f(2)=6$, if $k=\displaystyle\sum_{r=1}^{25}f(r)$, then find $\dfrac k{100}$

\qn{3} In a regular tetrahedron let $\theta$ be the angle between any edge and a face not containing the edge. If $\cos^2\theta=\dfrac ab$ where $a,b\in I^+$ also $a$ and $b$ are coprime, then find the value of $\dfrac{10a}b$

\qn{4} 5 girls and 10 boys sit at random in a row having 15 chairs numbered as 1 to 15. If the probability that the end seats are occupied by the girls and between any two girls odd number of boys take seat is $\dfrac n{3003}$, then find the value of $n$

\qn{5} A bag contains 5 white and 3 red balls. A ball is selected at random and discarded after noting its colour and, two balls of the other colour put in the bag. A ball now drawn from the bag is white and the probability that discarded ball was also white is $K$, then find $100K$.

\sect{SECTION-II(ii)}
\qn{6} If $C_1:x^2+y^2=(3+2\sqrt2)^2$ be a circle and $PA$ and $PB$ are pair of tangents on $C_1$ where $P$ is any point on the director circle of $C_1$, then the radius of smallest circle which touches $C_1$ externally and also the two tangents $PA$ and $PB$, is --

\qn{7} Origin $O$ is the centre of two concentric circles whose radii are $a$ \& $b$ respectively, $a<b$. A line $OPQ$ is drawn to cut the inner circle in $P$ \& the outer circle in $Q$. $PR$ is drawn parallel to the $y$-axis \& $QR$ is drawn parallel to the $x$-axis. The locus of $R$ is an ellipse touching the two circles and the focii of this ellipse lie on the inner circle. if eccentricity of that ellipse is $\lambda$, then find $\sqrt2\lambda$

\qn{8} Let $f(x)=\dfrac{\tan x}x$ and $\displaystyle\lim_{x\to0}\left([f(x)]+x^2\right)^{\frac1{\{f(x)\}}}=e^\lambda$, then find $\lambda$\\
{(where $[.]$ and $\{.\}$ denotes greatest integer and fractional part function respectively)}

\qn{9} If $f(x)=\begin{cases}\dfrac{\tan[x^2]\pi}{ax^2}+ax^3+b, & 0<x\le1\\[6pt] 2\cos\pi x+\tan^{-1}x, & 1<x\le2\end{cases}$ is differentiable in $(0,2]$, then $a=\dfrac1{k_1}$ and $b=\dfrac\pi4-\dfrac{26}{k_2}$.\\
Find $k_1^2+k_2^2$ {where $[.]$ denotes greatest integer function}.

\qn{10} Find the area enclosed by the solution set of $[x]\cdot[y]=2$. Where $[\cdot]$ represent greatest integer function of $x$.

\papersec{2}
\partsec{1}{PHYSICS}
\sect{SECTION-I}
\end{multicols}
\noindent\begin{minipage}{\textwidth}
\paraq{1}{2}
In column-1, the variation of potential energy (in J) as a function of $x$ is given for a particle. The particle moves under the effect of conservative force acting along the $x$-axis. Assume that no other force is acting on the particle. In column-2, locations on the $x$-axis are mentioned, and in column-3, various choices for the kinetic energy of particle are given.
\begin{center}\footnotesize
\begin{tabular}{|p{0.7cm}|p{5.2cm}|p{0.7cm}|p{2.2cm}|p{0.7cm}|p{2.4cm}|}
\hline
 & \centering Column--1\newline Potential energy as function of $x$ & \multicolumn{2}{c|}{\shortstack{Column--2\\Location on the $x$-axis}} & \multicolumn{2}{c|}{\shortstack{Column--3\\Kinetic energy of particle (in J)}}\\
\hline
(I) & $U(x)=\dfrac1{100}\left[x^3-6x^2+11x-6\right]$ & (i) & $x=0$ & (P) & KE $=0$\\
(II) & $U(x)=-2x^2+8x-6$ & (ii) & $x=1$ & (Q) & KE $=1$\\
(III) & $U(x)=\dfrac13\left[x^4-7x^3+17x^2-17x+6\right]$ & (iii) & $x=2$ & (R) & KE $=2$\\
(IV) & $U(x)=2x^2-4x$ & (iv) & $x=3$ & (S) & KE $=3$\\
\hline
\end{tabular}
\end{center}
\end{minipage}
\begin{multicols}{2}

\qn{1} In which of the following case the force directed towards $+$ve $x$-axis if mechanical energy of the particle is 3 J?
\begin{qoptions}
\opt{A}{(II) (ii) (S)}
\opt{B}{(I) (iv) (S)}
\opt{C}{(II) (iv) (S)}
\opt{D}{(II) (iii) (Q)}
\end{qoptions}

\qn{2} In which of the following the particle is in stable equilibrium if mechanical energy of the particle is 3 J?
\begin{qoptions}
\opt{A}{(III) (ii) (S)}
\opt{B}{(II) (iii) (Q)}
\opt{C}{(III) (i) (Q)}
\opt{D}{(II) (ii) (S)}
\end{qoptions}

\end{multicols}
\noindent\begin{minipage}{\textwidth}
\paraq{3}{4}
A person stands on a platform-and-pulley system, as shown in Column-I. The masses of the platform and person are $M$ and $m$ respectively. The rope and pulley is massless. The man can pull on the rope so as to produce one of the situations given in following question. In column-I $W_1$ is tension \& $W_2$ is normal reaction between the man and the platform. Match the possible situations.
\begin{center}\small
\begin{tabular}{|p{0.8cm}|p{3.6cm}|p{0.8cm}|p{3.8cm}|p{0.8cm}|p{3.8cm}|}
\hline
 & \centering Column-I & \multicolumn{2}{c|}{Column-II} & \multicolumn{2}{c|}{Column-III}\\
\hline
(I) & \qfigi{p2_phy_s1_para3_4_1}\rule{0pt}{35pt} & (i) & $W_2=(2m-M)\dfrac g3$ & (P) & $W_1=(m+M)\dfrac g3$\\
\hline
(II) & \qfigi{p2_phy_s1_para3_4_2}\rule{0pt}{45pt} & (ii) & $W_2=(2m-M)\dfrac g2$ & (Q) & $W_1=(m+M)\dfrac g2$\\
\hline
(III) & \qfigi{p2_phy_s1_para3_4_3}\rule{0pt}{42pt} & (iii) & $W_2=(m-M)\dfrac g4$ & (R) & $W_1=(m+M)g$\\
\hline
(IV) & \qfigi{p2_phy_s1_para3_4_4}\rule{0pt}{33pt} & (iv) & $W_2=(3m-M)\dfrac g4$ & (S) & $W_1=(m+M)\dfrac g4$\\
\hline
\end{tabular}
\end{center}
\end{minipage}
\begin{multicols}{2}

\qn{3} Platform is in equilibrium then :-
\begin{qoptions}
\opt{A}{(IV) (ii) (Q)}
\opt{B}{(IV) (iv) (R)}
\opt{C}{(IV) (iv) (Q)}
\opt{D}{(III) (ii) (Q)}
\end{qoptions}

\qn{4} Platform is accelerated by $\dfrac g2$ upward :-
\begin{qoptions}
\opt{A}{(I) (ii) (R)}
\opt{B}{(II) (iv) (Q)}
\opt{C}{(III) (i) (P)}
\opt{D}{(III) (ii) (R)}
\end{qoptions}

\paraq{5}{6}
In the given circuits, all the ammeters are ideal but voltmeters may not be ideal. In the below column-1 $A_1,A_2,A_3$ are used to represent the magnitude of current through the ammeters $A_1,A_2,A_3$. In Column-2 $V_1$ and $V_2$ are readings of voltmeters. $R_{eq}$ is equivalent resistance for different situations.
\qfig{p2_phy_s1_para5_6_1}
\begin{center}\footnotesize
\begin{tabular}{|p{0.6cm}|p{1.6cm}|p{0.5cm}|p{1.6cm}|p{0.5cm}|p{1.6cm}|}
\hline
 & \centering Column-1\newline $|A_2-A_1|$ & & \centering Column-2\newline $|V_2-V_1|$ & & \centering Column-3\newline $R_{eq}$\tabularnewline
\hline
(I) & $\dfrac{25}3$ & (i) & 30 & (P) & 5\\
(II) & 4 & (ii) & 3 & (Q) & 9\\
(III) & 3 & (iii) & 0 & (R) & 6\\
(IV) & 8 & (iv) & 4 & (S) & 4\\
\hline
\end{tabular}
\end{center}

\qn{5} If $R_1=3\,\Omega$, $R_2=6\,\Omega$. Resistance of voltmeters $V_1$ and $V_2$ are $6\,\Omega$ and $3\,\Omega$ respectively and $E=24$ V and both the switches are closed.
\begin{qoptions}
\opt{A}{(II) (ii) (R)}
\opt{B}{(IV) (iii) (R)}
\opt{C}{(II) (iii) (S)}
\opt{D}{(I) (iv) (Q)}
\end{qoptions}

\qn{6} $R_1=8\,\Omega$, $R_2=2\,\Omega$, $E=50$ V. Resistance of $V_1$ and $V_2$ is $12\,\Omega$ and $3\,\Omega$ respectively and both the switches are closed :-
\begin{qoptions}
\opt{A}{(I) (iv) (R)}
\opt{B}{(IV) (i) (Q)}
\opt{C}{(I) (i) (R)}
\opt{D}{(III) (iii) (S)}
\end{qoptions}

\end{multicols}
\noindent\begin{minipage}{\textwidth}
\paraq{7}{8}
Column-I represent a particle of mass $m$ which projected as represent in column-I. In column-II a box of mass $m$ have its intial velocity $u$ \& acceleration $a$ as shown. In column-III a block of mass $m$ is put on wedge.
\begin{center}\small
\begin{tabular}{|p{0.8cm}|p{3.0cm}|p{0.8cm}|p{3.6cm}|p{0.8cm}|p{4.0cm}|}
\hline
 & \centering Column-I & & \centering Column-II & & \centering Column-III\tabularnewline
\hline
(I) & \qfigi{p2_phy_s1_para7_8_1}\rule{0pt}{30pt} & (i) & \qfigi{p2_phy_s1_para7_8_5} & (P) & \qfigi{p2_phy_s1_para7_8_9}\\
\hline
(II) & \qfigi{p2_phy_s1_para7_8_2}\rule{0pt}{30pt} & (ii) & \qfigi{p2_phy_s1_para7_8_6} & (Q) & \qfigi{p2_phy_s1_para7_8_10}\\
\hline
(III) & \qfigi{p2_phy_s1_para7_8_3}\rule{0pt}{30pt} & (iii) & \qfigi{p2_phy_s1_para7_8_7} & (R) & \qfigi{p2_phy_s1_para7_8_11}\\
\hline
(IV) & \qfigi{p2_phy_s1_para7_8_4}\rule{0pt}{30pt} & (iv) & \qfigi{p2_phy_s1_para7_8_8} & (S) & \qfigi{p2_phy_s1_para7_8_12}\\
\hline
\end{tabular}
\end{center}
\end{minipage}
\begin{multicols}{2}

\qn{7} In which combination displacement of particle in vertical direction is zero?
\begin{qoptions}
\opt{A}{(II) (i) (Q)}
\opt{B}{(III) (iii) (R)}
\opt{C}{(II) (iv) (S)}
\opt{D}{(II) (i) (P)}
\end{qoptions}

\qn{8} In which combination speed of particle must increase just after the motion start?
\begin{qoptions}
\opt{A}{(I) (i) (Q)}
\opt{B}{(III) (ii) (P)}
\opt{C}{(II) (iii) (R)}
\opt{D}{(I) (iv) (R)}
\end{qoptions}

\sect{SECTION-II(i)}
\qn{1} Figure shows a soap film formed between two square figures made of a uniform wire. The bigger square is held while keeping it in a horizontal plane and the smaller square is slowly allowed to drop vertically. It reaches an equilibrium state after dropping a height $h$. Let surface tension of soap $=T$. Mass per unit length of wire $=\lambda$. Acceleration due to gravity $=g$. Given that $h=\dfrac{n\lambda ga}{\sqrt{4T^2-\lambda^2g^2}}$; find the value $n$.
\qfig{p2_phy_s2_q01_1}

\qn{2} A spherical ball of density $\rho$ and radius 0.003 m is dropped into a tube containing a viscous fluid. Viscosity of the fluid $=1.260$ N.m$^{-2}$ and its density $\rho_L=\rho/2=1260$ kg.m$^{-3}$. Find its terminal speed in cm/s. ($g=$ acceleration due to gravity $=10$ ms$^{-2}$)

\qn{3} A particle of charge $+q$ enters a region of uniform magnetic field $B$ (directed into the plane of paper) as shown in the figure. Particle is deflected by distance $d$ along $Y$ axis after travelling a distance of $a$ along $X$ axis. Find magnitude of linear momentum of particle in Newton-sec. (Given: $a=3$ m, $d=4$ m, $qB=0.8$ C-Tesla)
\qfig{p2_phy_s2_q03_1}

\qn{4} A neutral atom of atomic mass number 100 which is stationary at the origin in gravity free space emits an $\alpha$-particle ($A$) in $z$-direction. The product ion is $P$. A uniform magnetic field exists in the $x$-direction. Disregard the electro magnetic interaction between $A$ and $P$. If the angle of rotation of $A$ after which $A$ and $P$ will meet for the first time is $\theta$ radians, what is the value of $\theta$?

\qn{5} Shots are fired from the top of a tower and from bottom simultaneously at angles $45^\circ$ and $60^\circ$ respectively as shown. If horizontal distance of the point of collision is at a distance of $(\sqrt3+1)$ m from the tower, then find the height $h$ of tower (in cm).
\qfig{p2_phy_s2_q05_1}

\qn{6} In the circuit shown, a 30 V d.c. source gives a current 2.0 A as recorded in the ammeter $A$ and 30 V a.c. source of frequency 100 Hz gives a current 1.2 A. The inductive reactance (in ohm) is :-
\qfig{p2_phy_s2_q06_1}

\sect{SECTION-II(ii)}
\qn{7} A river is flowing with speed 3 km/hr west to east. A man swims with speed 5 km/hr in still water. Man is at south bank of the river. Match the List-I with direction of velocities of man w.r.t. ground in List-II.
\begin{center}\footnotesize
\begin{tabular}{|p{0.6cm}|p{4.6cm}|p{0.6cm}|p{1.2cm}|}
\hline
 & \centering List-I & \multicolumn{2}{c|}{List-II}\\
\hline
(P) & Man swims at an angle $127^\circ$ from river flow. & (1) & \qfigi{p2_phy_s3_q07_1}\rule{0pt}{16pt}\\
\hline
(Q) & Man swims right angle to river flow. & (2) & \qfigi{p2_phy_s3_q07_2}\rule{0pt}{16pt}\\
\hline
(R) & Man swims at an angle $143^\circ$ from river flow. & (3) & \qfigi{p2_phy_s3_q07_3}\rule{0pt}{16pt}\\
\hline
\end{tabular}
\end{center}

\qn{8} \qfig{p2_phy_s3_q08_1}
\begin{center}\footnotesize
\begin{tabular}{|p{0.6cm}|p{4.8cm}|p{0.6cm}|p{1.6cm}|}
\hline
 & \centering List-I & \multicolumn{2}{c|}{List-II}\\
\hline
(P) & Phase of current in $B$ leads voltage of $A$ by a phase $\theta$. & (1) & $0<\theta<\dfrac\pi2$\\
\hline
(Q) & Phase of voltage in $D$ leads current of $B$ by phase $\theta$. & (2) & $-\dfrac\pi2<\theta<0$\\
\hline
(R) & Phase of current in $C$ leads current of $A$ by phase $\theta$. & (3) & $\theta=0$\\
\hline
\end{tabular}
\end{center}

\qn{9} List-I contains some circuits and List-II contains some properties. In List-II current signifies instantaneous current and potential difference refers instantaneous potential difference.
\begin{center}\footnotesize
\begin{tabular}{|p{0.5cm}|p{4.3cm}|p{0.5cm}|p{1.5cm}|}
\hline
 & \centering List-I\newline Charge System & \multicolumn{2}{c|}{\shortstack{List-II\\Properties}}\\
\hline
(P) & \qfigi{p2_phy_s3_q09_1}\rule{0pt}{32pt}\par A closely separated conducting wire is wrapped over an insulating cylinder. Wire is connected to a resistance $R$ and an alternating source via a switch $S$, which is closed at $t=0$ & (1) & In branch $x$, current varies\\
\hline
(Q) & \qfigi{p2_phy_s3_q09_2}\rule{0pt}{27pt}\par A charged metallic sphere of radius $r$ is earthed with a resistance $R$ and a switch $S$, which is closed at $t=0$. & (2) & Between points $a$ and $b$ magnitude of potential difference increases always.\\
\hline
(R) & \qfigi{p2_phy_s3_q09_3}\rule{0pt}{43pt}\par A circuit (used to measure an unknown resistance) is in balanced conditions (Galvanometer reads zero) & (3) & In branch $x$, current is constant\\
\hline
\end{tabular}
\end{center}

\qn{10} Match the entries of list-I with the entries of list-II, which describes the angle of deviation of ray with angle of incidence.
\begin{center}\footnotesize
\begin{tabular}{|p{0.6cm}|p{2.9cm}|p{0.6cm}|p{2.4cm}|}
\hline
 & \centering List-I & \multicolumn{2}{c|}{List-II}\\
\hline
(P) & A ray is falling on a plane smooth mirror & (1) & \qfigi{p2_phy_s3_q10_1}\rule{0pt}{32pt}\\
\hline
(Q) & A ray is going from a rarer to denser medium & (2) & \qfigi{p2_phy_s3_q10_2}\rule{0pt}{32pt}\\
\hline
(R) & A ray is going from a denser to rarer medium & (3) & \qfigi{p2_phy_s3_q10_3}\rule{0pt}{32pt}\\
\hline
\end{tabular}
\end{center}

\partsec{2}{CHEMISTRY}
\sect{SECTION-I}
\end{multicols}
\noindent\begin{minipage}{\textwidth}
\paraq{1}{2}

\begin{center}\small
\begin{tabular}{|p{0.8cm}|p{6.2cm}|p{0.8cm}|p{4.0cm}|p{0.8cm}|p{3.0cm}|}
\hline
 & \centering Column-I & & \centering Column-II & & \centering Column-III\tabularnewline
\hline
(I) & \qfigi{p2_chem_s1_para1_2_1}\rule{0pt}{20pt} & (i) & Carbene intermediate & (a) & Haffmann product\\
\hline
(II) & \qfigi{p2_chem_s1_para1_2_2}\rule{0pt}{40pt} & (ii) & Carbocation intermediate & (b) & Aromatic compound\\
\hline
(III) & \qfigi{p2_chem_s1_para1_2_3}\rule{0pt}{32pt} & (iii) & Transition state & (c) & 2-methyl But-2-ene\\
\hline
(IV) & \qfigi{p2_chem_s1_para1_2_4}\rule{0pt}{36pt} & (iv) & Free radical intermediate & (d) & cis alkene\\
\hline
\end{tabular}
\end{center}
\end{minipage}
\begin{multicols}{2}

\qn{1} Correct set for following reaction
\begin{qoptions}
\opt{A}{I, ii, d}
\opt{B}{I, iii, a}
\opt{C}{I, i, c}
\opt{D}{I, iii, b}
\end{qoptions}

\qn{2} Correct set for following reaction :
\begin{qoptions}
\opt{A}{III, ii, c}
\opt{B}{IV, iii, d}
\opt{C}{II, iii, b}
\opt{D}{All}
\end{qoptions}

\end{multicols}
\noindent\begin{minipage}{\textwidth}
\paraq{3}{4}

\begin{center}\small
\begin{tabular}{|p{0.8cm}|p{5.0cm}|p{0.8cm}|p{3.8cm}|p{0.8cm}|p{2.8cm}|}
\hline
 & \centering Column-I\newline Compound & & \centering Column-II\newline Hybridisation and magnetic behaviour & & \centering Column-III\newline Isomerism\tabularnewline
\hline
(I) & [Cr(NH$_3$)$_4$Cl$_2$]Cl & (i) & Paramagnetic & (P) & Ionisation\\
(II) & [Ti(H$_2$O)$_5$Cl](NO$_3$)$_2$ & (ii) & Diamagnetic & (Q) & Geometrical\\
(III) & [Pt(en)(NH$_3$)Cl]NO$_3$ & (iii) & d$^2$sp$^3$ & (R) & Linkage\\
(IV) & [Co(NH$_3$)$_4$(NO$_2$)$_2$]NO$_3$ & (iv) & sp$^3$d$^2$ & (S) & Optical\\
\hline
\end{tabular}
\end{center}
\end{minipage}
\begin{multicols}{2}

\qn{3} Which is the only CORRECT combination?
\begin{qoptions}
\opt{A}{(I), (i), (Q)}
\opt{B}{(II), (iv), (P)}
\opt{C}{(I), (ii), (S)}
\opt{D}{(II), (iii), (Q)}
\end{qoptions}

\qn{4} Which is the only INCORRECT combination.
\begin{qoptions}
\opt{A}{(III), (ii), P}
\opt{B}{(IV), (ii), Q}
\opt{C}{(III), (ii), Q}
\opt{D}{(IV), (iii), Q}
\end{qoptions}

\end{multicols}
\noindent\begin{minipage}{\textwidth}
\paraq{5}{6}

\begin{center}\small
\begin{tabular}{|p{0.8cm}|p{2.6cm}|p{0.8cm}|p{4.8cm}|p{0.8cm}|p{5.2cm}|}
\hline
 & \centering Column-I\newline Reaction & & \centering Column-II\newline Change associated with reaction & & \centering Column-III\tabularnewline
\hline
(I) & N$_2\to$ N$_2^+$ & (i) & Electron is removed from gerade BMO & (P) & Magnetic behaviour of reactant and product is different\\
(II) & N$_2\to$ N$_2^-$ & (ii) & Electron is removed from ungerade BMO & (Q) & Magnetic behaviour of reactant and product is same\\
(III) & C$_2\to$ C$_2^+$ & (iii) & Electron is added to gerade ABMO & (R) & Bond order of product is more as compare to reactant\\
(IV) & O$_2\to$ O$_2^-$ & (iv) & Electron is added to ungerade ABMO & (S) & Bond order of reactant is more as compare to product\\
\hline
\end{tabular}
\end{center}
\end{minipage}
\begin{multicols}{2}

\qn{5} Select CORRECT combination.
\begin{qoptions}
\opt{A}{(I), (ii), P}
\opt{B}{(I), (i), Q}
\opt{C}{(II), (iii), P}
\opt{D}{(II), (iv), R}
\end{qoptions}

\qn{6} Select INCORRECT combination.
\begin{qoptions}
\opt{A}{(III), (ii), P}
\opt{B}{(III), (ii), S}
\opt{C}{(IV), (iii), S}
\opt{D}{(IV), (iv), Q}
\end{qoptions}

\end{multicols}
\noindent\begin{minipage}{\textwidth}
\paraq{7}{8}
$E^0_{\text{MnO}_4^-/\text{Mn}^{2+}}=1.51$ V; \ $E^0_{\text{Zn}^{2+}/\text{Zn}}=-0.76$ V\\
$E^0_{\text{Cl}_2/\text{Cl}^-}=1.36$ V; \ $E^0_{\text{Ag}^+/\text{Ag}}=0.8$ V\\
$E^0_{\text{Fe}^{3+}/\text{Fe}^{2+}}=0.77$\\
$K_{sp}(\text{AgCl})=10^{-10}$\\
$\dfrac{2.303RT}{F}=0.06$, \ ($\log2=0.3$)
\begin{center}\small
\begin{tabular}{|p{0.9cm}|p{8.8cm}|p{0.9cm}|p{1.8cm}|p{0.9cm}|p{1.8cm}|}
\hline
 & \centering Column-I & & \centering Column-II & & \centering Column-III\tabularnewline
\hline
(I) & Pt $|$ H$_{2(g)}$(2atm) $|$ H$^+$(0.1M) $\|$ Cl$^-$(0.1) $|$ Cl$_{2(g)}$ (2atm)$|$ Pt & (i) & 0.6 V & (P) & 0.48\\
(II) & Zn(s) $|$ Zn$^{+2}$(0.01M) $\|$ Ag$^+$(0.1M) $|$ Ag(s) & (ii) & 1.56 & (Q) & 1.56\\
(III) & Pt $|$ Fe$^{+2}$(1M), Fe$^{+3}$(1M) $\|$ MnO$_4^-$(1M), H$^+$(pH =2), Mn$^{+2}$(1M) $|$ Pt & (iii) & 0.74 V & (R) & 0.26\\
(IV) & Ag(s) $|$ AgCl(s) $|$ Cl$^-$(0.1M) $\|$ Ag$^+$ (0.1 M) $|$ Ag(s) & (iv) & 1.36 V & (S) & 1.498\\
\hline
\end{tabular}
\end{center}
\end{minipage}
\begin{multicols}{2}

\qn{7} Select the correct option
\begin{qoptions}
\opt{A}{I -- (i) -- R}
\opt{B}{II -- (iii) -- S}
\opt{C}{III -- (II) -- Q}
\opt{D}{IV -- (i) -- P}
\end{qoptions}

\qn{8} Lowest value of $E_{\text{cell}}$ is for Cell
\begin{qoptions}
\opt{A}{III}
\opt{B}{I}
\opt{C}{II}
\opt{D}{IV}
\end{qoptions}

\sect{SECTION-II(i)}
\qn{1} Identify total number of Essential amino acids\\
(a) Glycine \quad (b) Alanine \quad (c) Phenyl alanine \quad (d) Arginine\\
(e) Cysteine \quad (f) Histidine \quad (g) Proline \quad (h) Tryptophan

\qn{2} How many equivalents of base are consumed in the following reaction ?\\
$\underset{\text{(1 equivalent)}}{\text{CH}_3\text{CHO}}\xrightarrow{\text{Br}_2/\bar{\text{O}}\text{H}/\Delta}\text{HCBr}_3+\text{HCOO}^-$

\qn{3} Number of alkenes of formula C$_4$H$_8$ which give same product with HBr in presence as well as absence of peroxide are.

\qn{4} Ionic solid Na$^+$A$^-$ crystallise in rock salt type structure. 2.592 gm of ionic solid salt NaA dissolved in water to make 2 litre solution. The pH of this solution is 8. If distance between cation and anion is 250 pm. Calculate density of ionic solid (in gm/cm$^3$).\\
{(Given: p$K_w=13$, p$K_a$(HA)$=5$, $N_A=6\times10^{23}$)}

\qn{5} Molar heat capacities at constant pressure for $A$, $B$ and $C$ are 3, 1.5 and 2 J/K mol. The enthalpy of reaction and entropy of reaction, $A+2B\to3C$ are 20 kJ/mol and 15 J/K-mole at 300 K. Calculate $\Delta G$ (in kJ/mol) for the reaction.\\
$\dfrac12A+B\to\dfrac32C$

\qn{6} Find out the number of cation(s) which form(s) black ppt (soluble in hot and dil HNO$_3$) on passing H$_2$S gas into their salt solution.\\
Mg$^{2+}_{(aq.)}$, Cu$^{2+}_{(aq.)}$, Ba$^{2+}_{(aq.)}$, Fe$^{3+}_{(aq.)}$, Ag$^+_{(aq.)}$, Al$^{3+}_{(aq.)}$, Hg$^{2+}_{(aq.)}$, Pb$^{2+}_{(aq.)}$, Mn$^{2+}_{(aq.)}$

\sect{SECTION-II(ii)}
\qn{7} \begin{tabular}[t]{|p{0.65cm}|p{1.6cm}|p{0.65cm}|p{3.8cm}|}
\hline
 & \centering List-I & & \centering List-II\tabularnewline
\hline
(P) & Nylon-6 & (1) & Computer disc\\
(Q) & Bakelite & (2) & Tyre cords, Ropes\\
(R) & HDPE & (3) & Buckets, Dustbins \& Bottles, Pipes\\
\hline
\end{tabular}

\qn{8} \begin{tabular}[t]{|p{0.6cm}|p{2.0cm}|p{0.6cm}|p{2.0cm}|}
\hline
 & \centering List-I & & \centering List-II\tabularnewline
\hline
(P) & Paints & (1) & Gas\\
(Q) & Fog & (2) & Solid\\
(R) & Smoke & (3) & Liquid\\
\hline
\end{tabular}

\qn{9} \begin{tabular}[t]{|p{0.65cm}|p{2.4cm}|p{0.65cm}|p{2.8cm}|}
\hline
 & \centering List-I & & \centering List-II\tabularnewline
\hline
(P) & n-hexane + n-heptane & (1) & Maximum boiling azeotrope\\
(Q) & Acetone + Chloroform & (2) & Can be separated by distillation\\
(R) & Ethanol + water & (3) & Minimum boiling azeotrope\\
\hline
\end{tabular}

\qn{10} Match the catalyst with given process.
\begin{center}\footnotesize
\begin{tabular}{|p{0.6cm}|p{2.0cm}|p{0.6cm}|p{3.0cm}|}
\hline
 & \centering List-I\newline Catalyst & & \centering List-II\newline Process\tabularnewline
\hline
(P) & PdCl$_2$ & (1) & Ostwald process\\
(Q) & CuCl$_2$ & (2) & Deacon process\\
(R) & Pt/Rh & (3) & Wacker process\\
\hline
\end{tabular}
\end{center}

\partsec{3}{MATHEMATICS}
\sect{SECTION-I}
\paraq{1}{2}
Consider the differential equation $(1-x^2)\dfrac{dy}{dx}+xy=ax$, $x\ne1$.\\
The point through which the curve passes and value of $a$ is given in column 1.\\
Column 2 : represents the nature of the curve.\\
Column 3 : represents the various values related to curve.
\begin{center}\footnotesize
\begin{tabular}{|p{0.6cm}|p{1.6cm}|p{0.6cm}|p{1.4cm}|p{0.5cm}|p{2.2cm}|}
\hline
 & \centering Column 1 & & \centering Column 2 & & \centering Column 3\tabularnewline
\hline
(I) & $(0,3)$, $a=1$ & (i) & Circle & (P) & Length of latus rectum $=1$\\
(II) & $(0,2)$, $a=1$ & (ii) & Ellipse & (Q) & Center $(0,3)$\\
(III) & $(3,1)$, $a=2$ & (iii) & Parabola & (R) & Length of latus rectum $=\dfrac2{15}$\\
(IV) & $(4,2)$, $a=3$ & (iv) & Hyperbola & (S) & Eccentricity $=\dfrac3{2\sqrt2}$\\
\hline
\end{tabular}
\end{center}

\qn{1} Identify the correct combination ?
\begin{qoptions}
\opt{A}{(I) (iii) (P)}
\opt{B}{(II) (i) (Q)}
\opt{C}{(III) (iv) (S)}
\opt{D}{(IV) (iv) (P)}
\end{qoptions}

\qn{2} Identify the correct combination ?
\begin{qoptions}
\opt{A}{(I) (ii) (R)}
\opt{B}{(II) (ii) (Q)}
\opt{C}{(III) (iv) (Q)}
\opt{D}{(IV) (iv) (Q)}
\end{qoptions}

\paraq{3}{4}
Three alike dice are rolled simultaneously.\\
Column 1 : Sum of numbers appeared on the dice are given.\\
Column 2 : Number of ways in which the event can occur.\\
Column 3 : Probability of the event.
\begin{center}\footnotesize
\begin{tabular}{|p{0.6cm}|p{2.2cm}|p{0.6cm}|p{0.8cm}|p{0.5cm}|p{1.3cm}|}
\hline
 & \centering Column 1 & \multicolumn{2}{c|}{Column 2} & \multicolumn{2}{c|}{Column 3}\\
\hline
(I) & 6 & (i) & 3 & (P) & $\dfrac12$\\
(II) & 15 & (ii) & 10 & (Q) & $\dfrac5{108}$\\
(III) & less than 11 & (iii) & 28 & (R) & $\dfrac3{56}$\\
(IV) & greater than or equal to 11 & (iv) & 108 & (S) & $\dfrac5{28}$\\
\hline
\end{tabular}
\end{center}

\qn{3} Identify the correct combination ?
\begin{qoptions}
\opt{A}{(I) (ii) (Q)}
\opt{B}{(I) (i) (Q)}
\opt{C}{(II) (i) (R)}
\opt{D}{(IV) (iv) (P)}
\end{qoptions}

\qn{4} Identify the correct combination ?
\begin{qoptions}
\opt{A}{(I) (i) (R)}
\opt{B}{(II) (ii) (R)}
\opt{C}{(III) (iii) (P)}
\opt{D}{(IV) (iv) (P)}
\end{qoptions}

\paraq{5}{6}
Four mirrors are placed along the lines $y=0$, $x=1$, $y=6$ and $x=-10$, a ray of line is emerging from origin, it reflects from the mirrors and exits as soon as it reaches the intersection of any two mirrors.\\
Column 1 : Denotes slope of incident ray\\
Column 2 : Denotes number of reflections before it exits\\
Column 3 : Denotes length of path covered before it exists
\begin{center}\footnotesize
\begin{tabular}{|p{0.6cm}|p{1.4cm}|p{0.6cm}|p{1.4cm}|p{0.5cm}|p{1.8cm}|}
\hline
 & \centering Column 1 & & \centering Column 2 & & \centering Column 3\tabularnewline
\hline
(I) & $m=1$ & (i) & 8 & (P) & $10\sqrt{10}$\\
(II) & $m=2$ & (ii) & 6 & (Q) & $12\sqrt5$\\
(III) & $m=3$ & (iii) & 2 & (R) & $12\sqrt2$\\
(IV) & $m=-3$ & (iv) & 4 & (S) & $12\sqrt{10}$\\
\hline
\end{tabular}
\end{center}

\qn{5} Identify the correct combination ?
\begin{qoptions}
\opt{A}{(I) (iii) (S)}
\opt{B}{(II) (iv) (Q)}
\opt{C}{(III) (i) (P)}
\opt{D}{(IV) (iv) (Q)}
\end{qoptions}

\qn{6} Identify the correct combination ?
\begin{qoptions}
\opt{A}{(I) (iii) (Q)}
\opt{B}{(II) (iv) (P)}
\opt{C}{(III) (ii) (R)}
\opt{D}{(IV) (iv) (P)}
\end{qoptions}

\paraq{7}{8}
$P(a,b)$ is a point in 1$^{\text{st}}$ quadrant through which circles pass touching both the coordinate axes.\\
Column 1 : Coordinates of point $P$.\\
Column 2 : Sum of ordinates of centers of all such circles.\\
Column 3 : Distance between the centers of such circles.
\begin{center}\footnotesize
\begin{tabular}{|p{0.6cm}|p{1.4cm}|p{0.6cm}|p{1.4cm}|p{0.5cm}|p{1.8cm}|}
\hline
 & \centering Column 1 & & \centering Column 2 & & \centering Column 3\tabularnewline
\hline
(I) & $(2,3)$ & (i) & 8 & (P) & $4\sqrt3$\\
(II) & $(3,1)$ & (ii) & 12 & (Q) & 8\\
(III) & $(2,2)$ & (iii) & 10 & (R) & $8\sqrt2$\\
(IV) & $(4,2)$ & (iv) & 6 & (S) & $4\sqrt6$\\
\hline
\end{tabular}
\end{center}

\qn{7} Identify the correct combination ?
\begin{qoptions}
\opt{A}{(I) (iii) (P)}
\opt{B}{(II) (i) (R)}
\opt{C}{(III) (i) (Q)}
\opt{D}{(IV) (ii) (S)}
\end{qoptions}

\qn{8} Identify the correct combination ?
\begin{qoptions}
\opt{A}{(I) (iii) (S)}
\opt{B}{(II) (i) (Q)}
\opt{C}{(III) (i) (R)}
\opt{D}{(IV) (ii) (P)}
\end{qoptions}

\sect{SECTION-II(i)}
\qn{1} If $f(n)=\displaystyle\int_0^{\pi/2}\dfrac{\sin^2n\theta}{\sin^2\theta}\,d\theta$, $(n\in N)$, then the value of $\dfrac{f(15)+f(3)}{f(15)-f(9)}$ is

\qn{2} The value of $\displaystyle\lim_{n\to\infty}n^2\int_{-1/n}^{1/n}(2021\tan x+2022\cos x)\,|x|\,dx$

\qn{3} Given $a^2+b^2+c^2\le16$ and $S=3a+4b+10c$ where $a,b,c\in R$. If the maximum value of $S$ is $\sqrt N$, then the value of $N$ is equal to

\qn{4} Consider curves $C_1:|\arg(z-2)|=\dfrac\pi4$; $C_2:|\arg(z+2)|=\dfrac{3\pi}4$ and $C_3:|z-\alpha i|=4\sqrt2$, $\alpha\in R$ on the complex plane such that $C_3$ touches both $C_1$ and $C_2$, then sum of absolute value(s) of $\alpha$ is

\qn{5} Let $f$ be a differentiable function on $[1,5]$, $f(1)=0$, $|f'(x)|\le2|f(x)|$ on $[1,5]$, then the value of $\displaystyle\sum_{r=1}^5f(r)$ is

\qn{6} Let $A$ and $C$ are square matrices of same order, $\det(A)\ne0$ and $AC-CA=C^2$, then the number of values of $\det(C)$ is

\sect{SECTION-II(ii)}
\qn{7} Let the first term `$a$' of an infinite G.P. is the value of $x$ where function $f(x)=7+2x\log_e25-5^{x-1}-5^{2-x}$ has greatest value and common ratio `$r$' is equal to $\displaystyle\lim_{x\to0}\int_0^x\dfrac{t^2\,dt}{x^2\tan x}$. If $S$ is sum of infinite terms then
\begin{center}\footnotesize
\begin{tabular}{|p{0.6cm}|p{1.2cm}|p{0.6cm}|p{1.0cm}|}
\hline
 & \centering List-I & & \centering List-II\tabularnewline
\hline
(P) & $a$ & (1) & 4\\
(Q) & $\dfrac1r$ & (2) & 3\\
(R) & $S$ & (3) & 2\\
\hline
\end{tabular}
\end{center}

\qn{8} In a $\Delta PQR$, the side lengths $QR$, $RP$ and $PQ$ are consecutive positive integers in increasing order.
\begin{center}\footnotesize
\begin{tabular}{|p{0.6cm}|p{5.3cm}|p{0.6cm}|p{1.0cm}|}
\hline
 & \centering List-I & & \centering List-II\tabularnewline
\hline
(P) & If $z_1$, $z_2$ and $z_3$ are vertices of $P$, $Q$ and $R$ respectively in argand plane, such that $\left|\arg\left(\dfrac{z_1-z_3}{z_2-z_3}\right)\right|=2\left|\arg\left(\dfrac{z_3-z_1}{z_2-z_1}\right)\right|$, then biggest side of triangle is & (1) & 21\\
\hline
(Q) & Let $\vec p$, $\vec q$ and $\vec r$ be the position vectors of vertices $P$, $Q$ and $R$ respectively. If $(\vec r-\vec p)\cdot(\vec q-\vec r)=0$, then the value of $|\vec p\times\vec q+\vec q\times\vec r+\vec r\times\vec p|$ is equal to & (2) & 12\\
\hline
(R) & Let the equations $a_1x+b_1y+c_1=0$ and $a_2x+b_2y+c_2=0$ represent the lines $PQ$ and $PR$ respectively and $\left|\dfrac{a_1b_2-a_2b_1}{a_1a_2+b_1b_2}\right|=\dfrac43$, then the value of the semiperimeter of $\Delta PQR$ is & (3) & 6\\
\hline
\end{tabular}
\end{center}

\qn{9} Let $h(x)=4x^4-ax^3$ and $g(x)$ is a quadratic polynomial satisfying $g(0)+6=g'(0)-c=g''(0)+2b=0$ (where $a,b,c\in R$). If $y=h(x)$ and $y=g(x)$ intersect at 4 distinct points with non negative abscissae $x_i$, $i=1,2,3,4$ such that $\displaystyle\sum_{i=1}^4\dfrac i{x_i}=8$.
\begin{center}\footnotesize
\begin{tabular}{|p{0.6cm}|p{1.6cm}|p{0.6cm}|p{1.0cm}|}
\hline
 & \centering List-I & & \centering List-II\tabularnewline
\hline
(P) & $a$ & (1) & 5\\
(Q) & $c$ & (2) & 25\\
(R) & $\displaystyle\sum_{i=1}^4x_i$ & (3) & 20\\
\hline
\end{tabular}
\end{center}

\qn{10} \begin{center}\footnotesize
\begin{tabular}{|p{0.6cm}|p{5.3cm}|p{0.6cm}|p{1.0cm}|}
\hline
 & \centering List-I & & \centering List-II\tabularnewline
\hline
(P) & The function $f(x)$ is defined as $f(x)=\begin{cases}(-1)^{[x^2]} & \text{if } x<0\\[2pt] \displaystyle\lim_{n\to\infty}\dfrac1{1+x^n} & \text{if } x\ge0\end{cases}$ where $[x]$ denotes the greatest integer function. Number of integral values of $x$ in $[-2,9]$ for which $f(x)$ is continuous is & (1) & 9\\
\hline
(Q) & $f:R\to R$ and satisfies $f(1)=-1$, $f'(1)=3$. If $\displaystyle\int_1^2(2-x)f''(x)\,dx=7$, then $f(2)$ equals & (2) & 1\\
\hline
(R) & If $\displaystyle\int_a^b\left(4+3x-x^2\right)dx$ is maximum then $(a+b)$ is equal to \ $(b>a)$ & (3) & 3\\
\hline
\end{tabular}
\end{center}

\end{multicols}
\end{document}
